% Formula examples use the amsmath environments already loaded by Beamer.
\begin{frame}{行内公式与多行推导}
  对样本 $\mathbf{x}_i\in\mathbb{R}^{d}$，用线性模型预测标量目标 $y_i$：
  \begin{align*}
    \widehat{y}_i &= \mathbf{w}^{\mathsf T}\mathbf{x}_i+b, \\
    \mathcal{L}(\mathbf{w},b)
      &= \frac{1}{N}\sum_{i=1}^{N}(\widehat{y}_i-y_i)^2
         +\lambda\lVert\mathbf{w}\rVert_2^2, \\
    \nabla_{\mathbf{w}}\mathcal{L}
      &= \frac{2}{N}\sum_{i=1}^{N}(\widehat{y}_i-y_i)\mathbf{x}_i
         +2\lambda\mathbf{w}.
  \end{align*}
  \begin{notebox}[排版要点]
    行内公式用 \texttt{\$...\$}；多行公式用 \texttt{align*}，
    在等号前放置 \texttt{\&} 对齐，星号表示不显示编号。
  \end{notebox}
\end{frame}
\note{沿等号解释预测、损失与梯度，强调多行推导应对齐关键运算符。这是数学排版示例。}

\begin{frame}{矩阵、分段函数与公式说明}
  \begin{columns}[T,onlytextwidth]
    \begin{column}{.47\textwidth}
      \textbf{矩阵：批量预测}
      \[
        \mathbf{X}=\begin{bmatrix}
          \mathbf{x}_1^{\mathsf T}\\
          \vdots\\
          \mathbf{x}_N^{\mathsf T}
        \end{bmatrix},\qquad
        \widehat{\mathbf{y}}=\mathbf{X}\mathbf{w}+b\mathbf{1}.
      \]
      \texttt{bmatrix} 自动生成方括号；矩阵中的列用 \texttt{\&} 分隔。
    \end{column}
    \begin{column}{.47\textwidth}
      \textbf{分段函数：ReLU}
      \[
        f(z)=\begin{cases}
          z, & z\geq 0,\\
          0, & z<0.
        \end{cases}
      \]
      \texttt{cases} 排列函数值与条件；公式中的文字可用 \texttt{\string\text\{...\}}。
    \end{column}
  \end{columns}
  \medskip
  \begin{tipbox}[符号与说明]
    $N$ 是样本数，$d$ 是特征维数；$\mathbf{1}\in\mathbb{R}^{N}$ 为全一向量。
    公式与符号解释放在同一页，便于阅读。
  \end{tipbox}
\end{frame}
\note{左侧检查矩阵间距，右侧检查分段条件。下方强调符号定义应在使用前说明。}
